\chapter{Diseño y Desarrollo}
\label{sec:diseno_desarrollo}

En este capítulo se describe la implementación técnica del sistema desarrollado, así como la estructura del código empleada para la ejecución de los distintos experimentos. Cada sección del capítulo corresponde a un bloque funcional independiente, permitiendo una descripción clara y modular del proceso de desarrollo.

\section{Implementación del análisis jerárquico cuántico}
\label{sec:implementacion_analisis}
En esta sección se detalla el diseño y desarrollo del código utilizado para la preparación de los datos, la construcción de las características económicas basadas en momentos estadísticos, la implementación del modelo cuántico y la generación de las visualizaciones analizadas en el capítulo anterior.

\subsection{Arquitectura general del sistema}
\label{sec:arquitectura_sistema}

El sistema se ha diseñado siguiendo una estructura modular con el objetivo de facilitar la mantenibilidad, la escalabilidad y la reproducibilidad de los experimentos. El flujo general del proceso se divide en tres bloques principales:

\begin{itemize}
    \item Preparación y preprocesamiento de datos (extracción de indicadores estadísticos).
    \item Implementación del modelo de análisis (QHRP y distancias de Frobenius).
    \item Generación de resultados y visualizaciones (Heatmaps de distancias).
\end{itemize}

Cada uno de estos bloques se implementa mediante scripts y cuadernos de trabajo independientes, permitiendo ejecutar y modificar cada etapa sin afectar al resto del sistema.

\subsection{Preparación y almacenamiento de datos}
\label{sec:preparacion_almacenamiento}
Los datos financieros empleados corresponden a precios de cierre ajustados, obtenidos y procesados inicialmente en formato \texttt{CSV} para asegurar la compatibilidad con las librerías de análisis de datos. 

Durante la carga de los datos se realiza un proceso de limpieza consistente en la eliminación de observaciones incompletas (NaNs) y el tratamiento de valores atípicos, garantizando la coherencia temporal de las series utilizadas. A partir de los precios se calculan los retornos diarios, los cuales constituyen la entrada principal para la función de ingeniería de características.

\subsection{Construcción del tensor de características}
\label{sec:construccion_tensor}

Las características se almacenan en un tensor tridimensional con dimensiones:

\[
(N, T, P),
\]

donde $N$ representa el número de activos, $T$ el número de observaciones temporales y $P=6$ el número de características consideradas. 

Para cada activo se calculan ventanas móviles utilizando la librería \texttt{scipy.stats}. En lugar de indicadores técnicos tradicionales, se emplean los momentos de la distribución de retornos (media, varianza, asimetría y curtosis), permitiendo una representación matemática profunda del perfil de riesgo de cada activo.

\subsection{Estandarización y control numérico}
\label{sec:estandarizacion_control}

Con el objetivo de garantizar estabilidad numérica durante las etapas posteriores del análisis, especialmente antes del mapeo al circuito cuántico, las características se estandarizan de forma independiente. Este procedimiento evita la dominancia de magnitudes extremas (como la curtosis) y preserva la dinámica relativa de cada serie temporal.

\subsection{Implementación del modelo cuántico}
\label{sec:implementacion_modelo_cuantico}

El modelo cuántico se implementa mediante circuitos parametrizados. Cada observación multivariante de 6 dimensiones se transforma en un estado cuántico simulado mediante \texttt{qiskit}, a partir del cual se construyen matrices densidad individuales $\rho$.

Para cada activo, dichas matrices se agregan mediante un promedio temporal, obteniendo una matriz densidad representativa del comportamiento global del activo. Todas las operaciones se realizan mediante simuladores de estado (\texttt{statevector\_simulator}), garantizando un entorno reproducible.

\subsection{Cálculo de distancias y estructura jerárquica}
\label{sec:calculo_distancias_estructura}

Una vez obtenidas las matrices densidad promedio, se calcula la distancia de Frobenius entre todos los pares de activos. Estas distancias se almacenan en una matriz simétrica $D$.

El orden de los activos se determina mediante técnicas de clustering jerárquico utilizando el método de enlace simple (\textit{single linkage}), siguiendo estrictamente la metodología HRP original. Este procedimiento permite identificar la estructura de grupos de activos basándose en la proximidad de sus estados cuánticos.

\subsection{Generación de visualizaciones}
\label{sec:generacion_visualizaciones}

Las visualizaciones se generan mediante mapas de calor (\texttt{seaborn}) que representan las matrices de distancia cuántica. Se compara la matriz original (desordenada) con la matriz reordenada mediante la jerarquía cuántica, lo que permite observar visualmente la aparición de bloques cuasi-diagonales que indican agrupaciones de activos similares.

\section{Descarga y preparación de datos financieros}
\label{sec:descarga_preparacion_datos}

\subsection{Construcción de características económicas (Momentos Estadísticos)}
\label{sec:construccion_caracteristicas_economicas}

Para cada activo se construye un conjunto de seis características que describen la distribución de sus retornos, lo cual es fundamental para el enfoque de densidad cuántica:

\begin{itemize}
    \item \textbf{Retorno diario:} El valor puntual del retorno en el tiempo $t$.
    \item \textbf{Media móvil:} El retorno promedio en la ventana seleccionada.
    \item \textbf{Volatilidad:} La desviación estándar de los retornos (segundo momento).
    \item \textbf{Retorno acumulado:} El producto de los retornos en la ventana.
    \item \textbf{Asimetría (Skewness):} Medida de la falta de simetría en la distribución (tercer momento).
    \item \textbf{Curtosis (Kurtosis):} Medida del apuntamiento y riesgo de cola (cuarto momento).
\end{itemize}

Estas características se calculan mediante ventanas móviles, permitiendo capturar dinámicas estadísticas de corto y medio plazo que los modelos de correlación simple suelen ignorar.

\subsection{Persistencia de los datos}
\label{sec:persistencia_datos}

Finalmente, el tensor de características se almacena en formato binario (\texttt{.npy}), permitiendo una carga eficiente durante la ejecución de los modelos de análisis cuántico. Este diseño desacopla la etapa de preparación de datos del análisis posterior, facilitando la reproducibilidad.